\documentclass[10pt]{article}
\usepackage{compresr}
\cprSetHeader{Compresr SDK · Core}

\begin{document}

\cprTitleBlock{Compresr SDK}{Core Usage}{%
    Query-aware LLM context compression for Python and TypeScript. Drop
    \textbf{30–70\% of the tokens} you send to GPT, Claude, or Gemini
    without losing the answer your user actually asked for.}

% ===========================================================================
\section{Why Compresr}

Most LLM context is wasted. Retrieved snippets, tool outputs, and chat
history carry far more text than the model needs to answer the question
in front of it. Compresr reads your context \emph{together with the
query} and returns a shorter version that keeps the answer-bearing
content and trims the rest. Same passage, two different questions, two
different compressions — and 30–70\% fewer tokens billed by your model
provider.

One cloud API, two thin SDKs. Same on-the-wire contract. Call
\cprInlineCode{compress(...)} and ship.

% ===========================================================================
\section{Install \& quick start}

Python 3.9+ and Node 20+ are the only requirements. Grab an API key
(prefix \cprInlineCode{cmp\_}) from
\href{https://compresr.ai}{compresr.ai}.

\begin{shcode}
pip install compresr             # Python
npm install @compresr/sdk        # TypeScript
export COMPRESR_API_KEY=cmp_...
\end{shcode}

\textbf{Python}
\begin{pycode}
from compresr import CompressionClient, RateLimitError

client = CompressionClient()             # picks up COMPRESR_API_KEY
try:
    result = client.compress(
        context="Long passage to compress...",
        query="What is the main conclusion?",
        target_compression_ratio=0.5,    # drop ~50% of tokens
    )
except RateLimitError:
    result = None                        # back off, retry, or fall through

if result and result.success:
    print(result.data.tokens_saved, result.data.compressed_context)
else:
    print("compression skipped:", getattr(result, "error", "n/a"))
\end{pycode}

\textbf{TypeScript}
\begin{tscode}
import { CompressionClient } from '@compresr/sdk';

const client = new CompressionClient();   // picks up COMPRESR_API_KEY
const result = await client.compress({
  context: 'Long passage to compress...',
  query: 'What is the main conclusion?',
  targetCompressionRatio: 0.5,
});
if (result.success) {
  console.log(result.data.tokens_saved, result.data.compressed_context);
} else {
  console.warn('compression skipped:', result.error);
}
\end{tscode}

The response carries the compressed text, original and compressed token
counts, achieved ratio, and call duration — everything you need to log
savings and prove ROI to your finance team. Typed exceptions
(\cprInlineCode{RateLimitError}, \cprInlineCode{ModelNotFoundError},
\ldots) let you branch on intent, and transient \cprInlineCode{429}/\cprInlineCode{503}
responses retry automatically with jittered backoff.

% ===========================================================================
\section{Batch and async}

Compress up to \textbf{100 contexts per call} for RAG indexing, replay
jobs, or batch evals. Use one shared query, or pair each context with
its own. Python is sync-first with full async equivalents; TypeScript
is async-only.

\textbf{Python}
\begin{pycode}
# One shared query
batch = client.compress_batch(
    contexts=["Doc 1...", "Doc 2..."],
    queries="What is self-attention?",
)
# Or per-context queries
batch = client.compress_batch(inputs=[
    {"context": "Doc A...", "query": "What is X?"},
    {"context": "Doc B...", "query": "What is Y?"},
])

# Async on the same client
async with CompressionClient(api_key="cmp_...") as c:
    result = await c.compress_async(context="...", query="...")
\end{pycode}

\textbf{TypeScript}
\begin{tscode}
const batch = await client.compressBatch({
  contexts: ['Doc 1...', 'Doc 2...'],
  queries: 'What is self-attention?',
});
const batch2 = await client.compressBatch({ inputs: [
  { context: 'Doc A...', query: 'What is X?' },
  { context: 'Doc B...', query: 'What is Y?' },
]});
\end{tscode}

Streaming (\cprInlineCode{compress\_stream} / \cprInlineCode{compressStream})
is available on both SDKs when you want to render compressed output as
it arrives.

% ===========================================================================
\section{Key parameters}\label{sec:params}

Every compress call takes the same knobs. Defaults are filled by the
backend so you can start with two arguments and tune later.

\renewcommand{\arraystretch}{1.15}
\begin{longtable}{@{}p{4.0cm} p{2.2cm} p{7.6cm}@{}}
\toprule
\textbf{Parameter} & \textbf{Default} & \textbf{What it does} \\
\midrule
\endhead

\cprInlineCode{context} & required &
The text you want shorter. \\

\cprInlineCode{query} & \cprInlineCode{None} &
The question your LLM is answering. Drives what's kept vs. dropped. \\

\cprInlineCode{target\_compression\_ratio} \newline
(TS: \cprInlineCode{targetCompressionRatio}) &
backend default &
\cprInlineCode{0–1} = fraction of tokens to remove
(\cprInlineCode{0.5} drops $\sim$50\%). \cprInlineCode{>1} = Nx factor
(\cprInlineCode{4} = $\sim$4$\times$ smaller). \\

\cprInlineCode{coarse} & backend default &
\cprInlineCode{true} = paragraph-level (faster);
\cprInlineCode{false} = token-level (finer, slower). \\

\cprInlineCode{compression\_model\_name} \newline
(TS: \cprInlineCode{compressionModelName}) &
\cprInlineCode{"latte\_v1"} &
Which model to use (see below). \\

\cprInlineCode{heuristic\_chunking} \newline
(TS: \cprInlineCode{heuristicChunking}) &
\cprInlineCode{None} &
Structure-aware chunking — useful for Markdown, code, legal docs. \\

\cprInlineCode{disable\_placeholders} \newline
(TS: \cprInlineCode{disablePlaceholders}) &
\cprInlineCode{None} &
Drop placeholder markers (\cprInlineCode{<\dots>}) from removed spans. \\

\bottomrule
\end{longtable}

% ===========================================================================
\section{Choose a model}\label{sec:models}

\begin{itemize}\itemsep1pt
    \item \cprInlineCode{latte\_v1} — the default. Best quality on
          long-form prose, RAG snippets, and tool outputs.
    \item \cprInlineCode{latte\_v2} — faster variant.
          Lower latency for high-throughput agent loops and tight P99
          budgets.
    \item \textbf{Agentic variants} (\cprInlineCode{hcc\_}, \cprInlineCode{toc\_},
          \cprInlineCode{tdc\_} prefixed) target chat history, tool
          outputs, and tool discovery respectively — see the framework
          PDFs below.
\end{itemize}

\begin{cprNoteBox}
\textbf{Quick rule:} start with defaults (\cprInlineCode{latte\_v1},
ratio \cprInlineCode{0.5}). Drop to \cprInlineCode{latte\_v2} if you
need lower latency. Tune \cprInlineCode{target\_compression\_ratio}
based on your eval scores.
\end{cprNoteBox}

% ===========================================================================
\section{Pick the right framework integration}

Already on a framework? Skip the plumbing — each integration has its
own short PDF.

{\sloppy
\begin{itemize}\itemsep1pt
    \item \textbf{LangChain} — drop-in document compressor and
          retriever wrapper.
          (\cprInlineCode{python\_langchain},
          \cprInlineCode{typescript\_langchain})
    \item \textbf{LangGraph} — long-running agent memory compression
          for stores, checkpoints, and handoffs.
          (\cprInlineCode{python\_langgraph},
          \cprInlineCode{typescript\_langgraph})
    \item \textbf{LlamaIndex} — node postprocessor, chat memory, and
          query-engine wrappers.
          (\cprInlineCode{python\_llamaindex},
          \cprInlineCode{typescript\_llamaindex})
    \item \textbf{LiteLLM proxy} — server-side guardrail that
          compresses every request flowing through your proxy.
          (\cprInlineCode{python\_litellm})
\end{itemize}
\par}

Questions, pricing, or a custom deployment? Reach us at
\href{https://compresr.ai}{compresr.ai} or
\href{mailto:founders@compresr.ai}{founders@compresr.ai}.

\end{document}
